% ======================================================================
\section{The iteration in which the chain runs}
\label{sec:iteration}
% ======================================================================

All of the coupling lives inside MOCASSIN's MC iteration, the recursive
subroutine \rtn{iterateMC} (\src{iteration\_mod.f90:\allowbreak{}56--3345}), which
advances the iteration counter and calls itself until the stop test of
Section~\ref{sec:convergence} passes or the maximum number of iterations is
reached. The PDR part switches on at iteration \kw{pdrStartIter} (default 1,
\src{common\_mod.f90:\allowbreak{}893}); a larger value lets MOCASSIN first converge an
ionisation structure, a gas and a dust temperature without molecules. From
then on every iteration executes the steps of Table~\ref{tab:steps}, in that
order.

\begin{longtable}{@{}L{0.04\linewidth}L{0.25\linewidth}L{0.25\linewidth}L{0.38\linewidth}@{}}
\caption{One MC iteration with the PDR coupling active (3D build). Lines are in \src{iteration\_mod.f90} unless stated.}\label{tab:steps}\\
\toprule
\# & Step & Location & What it does \\
\midrule
\endfirsthead
\toprule
\# & Step & Location & What it does \\
\midrule
\endhead
1 & MC transfer & \rtn{energyPacketDriver}, :611--613 (sources), :655--660 (optional diffuse source) & propagates the energy packets of every source; an absorbed packet is re-emitted on the spot by the local gas or dust and followed until it escapes; accumulates $\Jste$ and $\Jdif$ (Section~\ref{sec:field}) \\
2 & estimator units & :924, :940 & multiplies by $10^{-9}$ (packet luminosity in $10^{36}$\,erg\,s$^{-1}$, cell volume in $10^{45}$\,cm$^3$, path lengths in cm) and divides by 8 for a \kw{symmetricXYZ} octant \\
3 & dust per cell & \rtn{calibratePDRdust} :986 (first time), \rtn{dustPerCell} :993 & $\RFUV$, $\AV/N_{\mathrm H}$, grain area per H, model-wide fallbacks (Section~\ref{sec:dust}) \\
4 & columns and $\AV$ & \rtn{computeAVdust} :995 & \rtn{snapDiffuseSources} (snapshot for the diffuse walk), \rtn{computeAVradial} ($\AV$ and source-ward columns) (Section~\ref{sec:geometry}) \\
5 & field summary & \rtn{mocassinToPDR} :1023 & $\chi$ (local) and $\chi_0$ (geometric, unattenuated) per cell \\
6 & seed (first PDR iteration only) & :1027--1029, :1037 & \rtn{seedMolecules}, then steps 4--5 again on the seeded abundances, then \rtn{pdrToMocassin} (Section~\ref{sec:seed}) \\
7 & \Htwo{} level model & \rtn{h2LevelModel(grid,1)} :1043 & one sweep over all cells on this iteration's field (Section~\ref{sec:h2model}) \\
8 & MOCASSIN cell update & \rtn{updateCell} :1261 & ionisation and thermal balance of every cell; stores the atomic photoionisation and recombination rates the chemistry will use \\
9 & chemistry & \rtn{pdrStep} :1986, cells dealt round-robin to the ranks (:1980--1986), then \verb|MPI_Allreduce| of abundances, heating, cooling and rates & photorates, steady-state solve, heating and cooling for the thermal balance (Sections~\ref{sec:photorates}--\ref{sec:chemistry}) \\
10 & handback & \rtn{pdrToMocassin} :2154 & fraction of each element not in molecules (\rtn{elemAvail}) \\
11 & second ionisation solve & \rtn{ionSecondSolve} called at :2168 (routine :3361--3434); its control pass, before the chemistry, at :1274 & \rtn{ionBalance} solved again on the channels and molecular content this iteration's chemistry left, then put back (\dlog{D-404}); the result is set aside for step 15 \\
12 & damping, hydrogen budget & :2424 (only if \kw{iterDamp}$>0$; default 0), \rtn{rescaleIonToAvail} :2474 & ion stages rescaled to \rtn{elemAvail} so that $x(\mathrm{H^0})+x(\mathrm{H^+})+2x(\mathrm{H_2})=1$ \\
13 & stop test & :3243 & integrated-quantity test (\kw{convTotals}, Section~\ref{sec:convergence}) \\
14 & next iteration & recursive call of \rtn{iterateMC} & the opacity of the next transfer is built from this iteration's ionisation, dust and $n(\mathrm{H_2})$ \\
15 & finish & \rtn{ionKeepSecond} :3249/:3282 (routine :3455--3521), \rtn{h2LevelModel(grid,\allowbreak{}12)} :3254/:3283, outputs & the last second solve replaces the first in the written model; level model swept to convergence on the final model; \src{h2Cells.out}, \src{lineFluxH2.out}, \src{h2Pop.bin} \\
\bottomrule
\end{longtable}

\subsection{What each step reads, and from when}
\label{sec:lags}

The coupling is a fixed-point iteration in which some inputs are one
iteration old. Within iteration $n$:

\begin{itemize}
\item the \textbf{radiation field} used by steps 3--9 is iteration $n$'s,
  transported through the opacity built from iteration $n-1$'s ionisation,
  dust and $n(\mathrm{H_2})$;
\item the \textbf{shielding columns} (steps 4 and 7) are integrated over the
  abundances left by iteration $n-1$'s chemistry. The snapshot is taken
  before any cell of iteration $n$ is solved, so the columns do not depend
  on the rank count or on the order in which the cells are solved;
\item the \textbf{level model} (step 7) runs before \rtn{updateCell}: it sees
  iteration $n-1$'s $T_e$, collider abundances and $n(\mathrm{H_2})$, and
  iteration $n$'s field; the other cells' level populations in its columns
  are those of its previous sweep (a Jacobi iteration);
\item \textbf{MOCASSIN's ionisation balance} (step 8) takes the chemistry's
  channels to and from molecules, and normalises the ion stages to the
  molecular content, of iteration $n-1$ (\dlog{L-47}); step 11 measures what
  the delay costs and, at the end of the run, removes it from the written
  ions;
\item the \textbf{chemistry} (step 9) sees iteration $n$'s $T_e$, $n_e$,
  ionisation and atomic rates (step 8), iteration $n$'s field, the step-7
  Solomon rates and the columns of step 4;
\item the \textbf{coolant line escape} in step 9 integrates the optical depth
  over the level populations of the other cells as stored on the rank; these
  are iteration $n-1$'s except for cells solved earlier in the same sweep on
  the same rank, so this one input does depend weakly on the order of the
  sweep and the rank count;
\item the \textbf{heating terms} that MOCASSIN takes from the chemistry and
  the level model (\rtn{pdrHeatXtra}) and its \textbf{CO cooling table} enter
  step 8 as step 9 of iteration $n-1$ left them, and the heating terms are
  held fixed during the temperature search (Section~\ref{sec:thermal}). The
  cosmic-ray, turbulent and chemical terms were evaluated at $T_e^{(n-1)}$; the
  \Htwo{} terms come from the level model of iteration $n-1$, which ran at
  $T_e^{(n-2)}$.
\end{itemize}

Inside one iteration the chemistry is solved once per cell: every input of
the solve is frozen during the iteration, so repeating it would reproduce the
same answer (\dlog{D-354}, which retired the inner coupling loop after
finding that a column rescale between passes was the only thing that changed
and that it oscillated). The damping across iterations (\kw{iterDamp}) is off
by default (\dlog{D-406}): with it, the written model was a blend of solves
that trailed the undamped one by three to six iterations and, where a cell
has two solutions, wrote states no solve gave (\dlog{L-49}).

At a fixed point of the iteration every lagged input equals its current
value, so the lags determine the path to the solution, not the solution. Two
things qualify this. First, the MC field is redrawn at every iteration, so
the fixed point is reached only in a statistical sense: per-cell quantities
fluctuate about it with the MC noise of Section~\ref{sec:mcnoise}, and
because the photorates are non-linear in the field and the shielding is
non-linear in the columns, the noise can bias cell means (\dlog{L-26},
\dlog{D-408}). Second, the coupled system has cells with two stable states
(the carbon ionisation turn, \dlog{D-405}), in which the lagged inputs decide
which state a cell is in at a given iteration.
